% !TeX root = ../main.tex
\chapter{Introduction}
\chapterauthor{Aden Chen and Shravan Hari}

The aim of this series of lectures is to expose students to a wide range of topics in Microeconomic Theory.
We hope to deliver economic insights with minimal mathematical formalism, though this means that we sometimes have to sacrifice full generality to present a cleaner story.
% At the same time, we hope to preserve mathematical rigor as much as possible.
% In instances where we have to present only a sketch proof, we will clearly indicate so, or simply demote it from the glory of being in a proof environment.

This set of notes are written under the assumption that the reader has some familiarity with basic calculus, linear algebra, and probability.
Any more advanced mathematical tools required will be introduced along the way.
Some exposure to game theory is also assumed, though we make minimal use of its formal language.

For a taste of the topics, here are some of the questions we will analyze: 
\begin{itemize}
  \item Disagreement in beliefs can be caused by the parties having different information.
    Thus the disagreement itself should cause the parties to revise their own beliefs.
    To what extent can such revision happen?
    Is it possible to ``agree to disagree``?
  \item The same good can be auctioned off in multiple ways.
    One can, for instance, make the high bidder pay her bid (a \vocab{first price auction}), or ask each bidder to pay her own bid (an \vocab{all pay auction}).
    Which auction format maximizes profit?
    It turns out that under some conditions, all sensible auctions are the same.
  \item How do the different electoral systems compare, and how should one choose which to use?
  \item To what extent will an agent communicate her information to a decision maker?
    There will of course be full information revelation when the parties have perfectly aligned interests and no revelation when they have fully opposing interests (a \vocab{zero-sum game}), but what happens when the parties have imperfectly aligned preferences?
  \item To what extent can agents' information change the outcome of a game?
  \item How should a firm optimally design a contract to induce a worker to exert the right amount of effort?
    What if the firm cannot observe the worker's action (but only their output)?
    What if the firm doesn't even know what actions the worker has?
\end{itemize}

And here are some buzzwords: game theory, mechanism design, information economics, auction theory, contract theory, information design, strategic communication, robust mechanism design, global games.


